\centerline{\bf \Huge Вторая МатДрака }

\problem{\textbf{1}} Замените звёздочки цифрами так, чтобы равенство стало верным
и все семь цифр были различными: *** + ** = 1056.


\problem{\textbf{2}} В саду у Леонардо и Микеланджело росло 2026 розовых кустов. Микеланджело
полил половину всех кустов, и Леонардо полил половину всех кустов. При этом оказалось, что ровно
19 кустов, самых красивых, были политы и Леонардо, и Микеланджело. Сколько розовых кустов остались
не политыми?

\problem{\textbf{3}} В шестизначном числе первую и последнюю цифру заменили
звёздочками: *2015*. Известно, что число делится на 72. Восстановите число.


\problem{\textbf{4}} Нарисуйте шесть лучей так, чтобы они пересекались ровно в четырех точках, по три луча в каждой точке. Отметьте начала лучей жирными точками.


\problem{\textbf{5}} Решите ребус: БАО $\times$ БА $\times$ Б = 2002.


\problem{\textbf{6}} Среди 35 учеников 5 «А» класса не любят есть конфеты 13, торты —
12, пряники — 9 учеников. Кроме того, не любят есть конфеты и торты 3, пряники и торты — 6,
конфеты и пряники — 5 учеников. Наконец, не любят есть конфеты, торты и пряники 2 ученика.
Сколько в классе ребят, которые любят есть конфеты, торты и пряники?


\problem{\textbf{7}} Среди некоторых 13 последовательных натуральных чисел 7 чётных и 5 делятся на три. Сколько среди них чисел, которые делятся на 6?


\problem{\textbf{8}} В группе 17 человек знают английский язык, 14 человек знают
китайский язык, 20 человек знают арабский язык и 19 человек знают польский язык. При этом
34 человека в группе знают ровно один язык из перечисленных, а остальные — ровно два языка
из перечисленных. Сколько человек в группе?


\problem{\textbf{9}} Сколько натуральных чисел от 1 до 2015 включительно имеют сумму цифр, которая делится на 5?


\problem{\textbf{10}} Решите ребус: ABCDEF $\times$ 3 = BCDEFA.


\problem{\textbf{11}} Начертите два четырехугольника с вершинами в узлах сетки, из которых можно сложить и треугольник, и четырехугольник, и пятиугольник. Покажите, как это можно сделать.


\newpage
\hspace{6mm}

\problem{\textbf{12}} На доске записаны все числа от 1 до 300 включительно. Леди Баг раскрашивает в оранжевый цвет только те числа, у которых общий делитель с числом 99 больше 1.
Сколько всего чисел станут оранжевыми?


\problem{\textbf{13}} 
Решите ребус: MSU + MSU + MSU + MSU + OLYMP + OLYMP = MOSCOW.


\problem{\textbf{14}} Из натурального числа вычли сумму его цифр, из
полученного числа снова вычли сумму его (полученного числа) цифр и т.д. После одиннадцати
таких вычитаний получился нуль. С какого числа начинали? Перечислите все варианты!




\problem{\textbf{15}} 
Решите ребус: 250 $\times$ ЛЕТ + МГУ = 2005 $\times$ ГОД.

